\section{Алгоритмические языки. Представление о языках программирования высокого уровня. Пакеты прикладных программ}

\subsection{Алгоритм и алгоритмический язык}

\textbf{Алгоритм} --- точно заданная конечная процедура, преобразующая допустимые исходные данные в результат. Обычно выделяют дискретность, определённость действий, результативность и применимость к некоторому классу однотипных входов.

Алгоритм можно задавать формулами, словесным описанием, блок-схемой, псевдокодом или программой.

\textbf{Алгоритмический язык} --- формальный язык для однозначной записи алгоритмов. Он задаёт алфавит, синтаксические правила, представление данных и управляющие конструкции. \textbf{Язык программирования} --- алгоритмический язык с формально определённой семантикой и средствами трансляции/выполнения на вычислительной системе.

Для структурного программирования основными конструкциями являются следование, ветвление и цикл.

\subsection{Языки высокого уровня}

Высокоуровневый язык скрывает большую часть деталей машинных команд и предоставляет абстракции данных и алгоритмов. Примеры: Fortran, C/C++, Java, Python, C\#.

\textbf{Синтаксис} определяет форму допустимых конструкций; \textbf{семантика} --- их смысл. Тип данных задаёт множество значений и допустимых операций. Типизация может быть статической или динамической.

Основные парадигмы:
\begin{itemize}
\item процедурная/императивная;
\item объектно-ориентированная;
\item функциональная;
\item логическая/декларативная.
\end{itemize}
Большинство современных языков мультипарадигменны.

\subsection{Трансляция и выполнение}

\textbf{Компилятор} преобразует исходную программу в машинный код или промежуточное представление. Типичные стадии:
\[
\text{лексический}
\rightarrow
\text{синтаксический}
\rightarrow
\text{семантический анализ}
\rightarrow
\text{оптимизация}
\rightarrow
\text{генерация кода}.
\]
\textbf{Интерпретатор} организует выполнение программных конструкций без традиционной предварительной генерации отдельного машинного файла. Реальные реализации могут сочетать интерпретацию, байткод и JIT-компиляцию.

Высокоуровневое программирование опирается на стандартные библиотеки и внешние программные модули. Поэтому практическая ценность языка определяется не только его синтаксисом, но и экосистемой численных, графических, сетевых и других библиотек.

\subsection{Пакеты прикладных программ}

\textbf{Пакет прикладных программ} --- набор взаимосвязанных программных средств для решения определённого класса прикладных задач. В математическом моделировании типичная структура инженерного комплекса:
\[
\boxed{\text{препроцессор}\rightarrow\text{решатель}\rightarrow\text{постпроцессор}}.
\]

Препроцессор задаёт геометрию, параметры, начальные и граничные условия и сетку; решатель реализует численный алгоритм; постпроцессор анализирует и визуализирует результат. Пакет может дополнительно включать базы материалов, оптимизатор, API, язык сценариев и средства параллельных вычислений.

Следует различать язык программирования и ППП: язык является средством создания алгоритмов и программ общего назначения, а пакет предоставляет готовые модели/алгоритмы и инфраструктуру для заданного класса задач. Многие пакеты при этом позволяют расширять возможности пользовательским кодом.

\subsection{Выбор инструментов}

Для численного проекта язык выбирают по производительности, переносимости, системе типов, доступным библиотекам и возможностям параллельного выполнения. ППП выбирают по поддерживаемым моделям, численным методам, сеткопостроению, масштабируемости, средствам автоматизации и верификации.

\subsection{Что важно сказать на экзамене}

Нужно последовательно определить
\[
\boxed{\text{алгоритм}\rightarrow\text{алгоритмический язык}\rightarrow\text{язык высокого уровня}\rightarrow\text{ППП}}.
\]
Затем объяснить синтаксис/семантику, компиляцию/интерпретацию и привести структуру прикладного программного комплекса.

\newpage
